\section{Global Performance}
\subsection{Global Performance}
L'entraînement initial sur le dataset synthétique donne les résultats suivants à convergence : 
\begin{itemize}
    \item MSE Loss : 0.0006 (entraînement et validation)
    \item MAE : 0.016 (entraînement et validation)
\end{itemize}
Les courbes montrent une convergence rapide en 2-3 epochs sans overfitting. 
La val loss légèrement inférieure à la train loss confirme que le dataset 
synthétique est homogène et peu varié. \\

\subsection{Fine-tuning (dataset réel)}
Le fine-tuning sur le dataset réel donne les résultats suivants à convergence :
\begin{itemize}
    \item MSE Loss : 0.017 (validation)
    \item MAE : 0.098 (validation)
\end{itemize}
Les courbes de train et validation descendent conjointement sans divergence, 
indiquant une bonne généralisation sans overfitting. 

\subsection{Analyse comparative}
La MAE est environ 6 fois plus élevée après fine-tuning sur données réelles 
(0.098) qu'après entraînement sur données synthétiques (0.016). 
Cette différence s'explique par la plus grande variabilité des données réelles : 
bruit, variations d'éclairage et calibration spécifique du robot.

Cependant, les métriques numériques ne sont pas le seul critère de succès. 
Le vrai test est le comportement du robot sur piste : le modèle entraîné 
uniquement sur données synthétiques sortait régulièrement du circuit, 
tandis que le fine-tuning vise à corriger ce comportement en adaptant 
le modèle à la réalité physique du robot. \\

\subsection{Tests sur le robot}
Par manque de temps, nous n'avons pas pu tester le modèle final sur le robot.
